\chapter{Implementation}

\section{Development Environment}

The project was developed and tested in the following environment:

\begin{table}[H]
\centering
\caption{Development Environment Specifications}
\label{tab:devenv}
\begin{tabular}{|l|l|}
\hline
\textbf{Component} & \textbf{Details}\\
\hline
Operating System & Windows 11 / Ubuntu 22.04 LTS\\
\hline
Python Version & Python 3.10.x\\
\hline
Virtual Environment & venv (standard Python)\\
\hline
IDE & Visual Studio Code\\
\hline
Key Libraries & See Table~\ref{tab:dependencies}\\
\hline
Version Control & Git + GitHub\\
\hline
\end{tabular}
\end{table}

\begin{table}[H]
\centering
\caption{Python Library Dependencies}
\label{tab:dependencies}
\begin{tabular}{|l|l|p{5cm}|}
\hline
\textbf{Library} & \textbf{Version} & \textbf{Purpose}\\
\hline
pandas & $\geq$2.0 & Data loading, manipulation, rolling window operations\\
\hline
numpy & $\geq$1.24 & Numerical operations, threshold scanning\\
\hline
scikit-learn & $\geq$1.3 & Train/test split, metrics, confusion matrix\\
\hline
xgboost & $\geq$2.0 & Gradient boosting classifier\\
\hline
networkx & $\geq$3.0 & Transaction graph construction and centrality\\
\hline
imbalanced-learn & $\geq$0.11 & SMOTE and resampling utilities\\
\hline
fastapi & $\geq$0.100 & REST API framework\\
\hline
uvicorn & $\geq$0.23 & ASGI server for FastAPI\\
\hline
streamlit & $\geq$1.25 & Interactive dashboard\\
\hline
plotly & $\geq$5.15 & Gauge charts and visualizations in dashboard\\
\hline
shap & $\geq$0.42 & SHAP explainability (planned for next version)\\
\hline
requests & $\geq$2.31 & HTTP client for API testing\\
\hline
\end{tabular}
\end{table}

\section{Installation and Setup}

\subsection{Cloning the Repository}

\begin{lstlisting}[language=bash, caption={Repository setup commands}]
git clone https://github.com/smrutiranjanbhuyan/upi-guard.git
cd upi-guard
\end{lstlisting}

\subsection{Creating the Virtual Environment}

\begin{lstlisting}[language=bash, caption={Virtual environment setup}]
# Create virtual environment
python -m venv venv

# Activate on Windows
.\venv\Scripts\Activate.ps1

# Activate on Linux/macOS
source venv/bin/activate
\end{lstlisting}

\subsection{Installing Dependencies}

\begin{lstlisting}[language=bash, caption={Installing project dependencies}]
pip install -r requirements.txt
\end{lstlisting}

\section{Implementation of the Preprocessing Module}

\subsection{Complete build\_features() Function}

The complete preprocessing module (\texttt{src/preprocess.py}) is the backbone of the system. Below is an annotated walkthrough of its implementation.

\begin{lstlisting}[caption={Stage 1 -- Data cleaning and timestamp parsing}, label={lst:stage1}]
def build_features(df):
    df = df.copy()  # Avoid mutating the original DataFrame

    # Parse timestamps from string to datetime
    df["timestamp"] = pd.to_datetime(df["timestamp"],
                                      errors="coerce")
    # Drop records where timestamp parsing failed
    df = df.dropna(subset=["timestamp"])

    # Sort by sender and time -- essential for rolling windows
    df = df.sort_values(["sender_id", "timestamp"])
\end{lstlisting}

The \texttt{df.copy()} call at the start is important: it prevents the function from modifying the original DataFrame passed by the caller, which could cause subtle bugs if the caller reuses that DataFrame after the function returns.

\begin{lstlisting}[caption={Stage 2 -- Temporal feature extraction}, label={lst:stage2}]
    df["hour"] = df["timestamp"].dt.hour
    df["day_of_week"] = df["timestamp"].dt.dayofweek
    df["is_weekend"] = (df["day_of_week"] >= 5).astype(int)
    df["month"] = df["timestamp"].dt.month

    # Night transactions: 12:00 AM to 5:00 AM
    df["is_night"] = (
        (df["hour"] >= 0) & (df["hour"] <= 5)
    ).astype(int)

    # Salary week: first 7 days of the month
    df["is_salary_week"] = (
        df["timestamp"].dt.day <= 7
    ).astype(int)
\end{lstlisting}

\begin{lstlisting}[caption={Stage 3 -- Transaction velocity with rolling window}, label={lst:stage3}]
    # Must set timestamp as index for time-based rolling
    df = df.set_index("timestamp")

    df["txn_velocity_1h"] = (
        df.groupby("sender_id")["amount"]
        .rolling("1h")         # 1-hour time window
        .count()               # Count transactions in window
        .reset_index(level=0, drop=True)
    )

    df = df.reset_index()
    df["txn_velocity_1h"] = df["txn_velocity_1h"].fillna(1)
\end{lstlisting}

The \texttt{reset\_index(level=0, drop=True)} call in the rolling computation is a pandas idiom that removes the grouped key (\texttt{sender\_id}) from the multi-level index returned by \texttt{groupby().rolling()}, leaving only the timestamp index to be re-merged with the original DataFrame.

\begin{lstlisting}[caption={Stage 4 -- Behavioral profiling per sender}, label={lst:stage4}]
    # Average transaction amount for this sender
    df["sender_mean_amt"] = (
        df.groupby("sender_id")["amount"].transform("mean")
    )

    # Standard deviation of amounts
    df["sender_std_amt"] = (
        df.groupby("sender_id")["amount"].transform("std")
    )

    # Total transactions from this sender
    df["sender_txn_count"] = (
        df.groupby("sender_id")["amount"].transform("count")
    )

    # Number of unique counterparties
    df["unique_receivers"] = (
        df.groupby("sender_id")["receiver_id"]
        .transform("nunique")
    )

    # Z-score: how unusual is this amount for this sender?
    df["amount_zscore"] = (
        (df["amount"] - df["sender_mean_amt"]) /
        (df["sender_std_amt"] + 1e-5)
    )
\end{lstlisting}

\begin{lstlisting}[caption={Stage 5 -- Graph construction and centrality features}, label={lst:stage5}]
    G = nx.from_pandas_edgelist(
        df,
        source="sender_id",
        target="receiver_id",
        create_using=nx.DiGraph()   # Directed graph
    )

    # PageRank: measure of node importance/centrality
    pagerank = nx.pagerank(G)

    # Degree centrality: fraction of nodes connected to
    degree = nx.degree_centrality(G)

    # Map graph metrics back to each row
    df["sender_pagerank"] = df["sender_id"].map(pagerank)
    df["sender_degree"]   = df["sender_id"].map(degree)

    # Fill for any senders not found in graph (shouldn't
    # happen, but defensive programming)
    df["sender_pagerank"] = df["sender_pagerank"].fillna(0)
    df["sender_degree"]   = df["sender_degree"].fillna(0)
\end{lstlisting}

\begin{lstlisting}[caption={Stages 6-8 -- Cross-state flag, encoding, cleanup}, label={lst:stages6_8}]
    # Cross-state transaction flag
    if ("sender_state" in df.columns and
            "receiver_state" in df.columns):
        df["cross_state"] = (
            df["sender_state"] != df["receiver_state"]
        ).astype(int)

    # One-hot encode categorical features
    categorical_cols = [
        "transaction_type", "merchant_category",
        "sender_state",     "receiver_state",
        "sender_bank",      "receiver_bank",
        "device_type",      "network_type"
    ]

    existing_cols = [c for c in categorical_cols
                     if c in df.columns]

    df = pd.get_dummies(df, columns=existing_cols,
                        drop_first=True)

    # Final cleanup: replace NaN with 0
    df = df.fillna(0)

    return df
\end{lstlisting}

\section{Implementation of the Training Module}

\subsection{Data Loading and Feature Building}

\begin{lstlisting}[caption={Data loading and feature building in train.py}, label={lst:trainload}]
import pandas as pd
import numpy as np
import pickle
import xgboost as xgb
from sklearn.model_selection import train_test_split
from sklearn.metrics import (classification_report,
    roc_auc_score, confusion_matrix,
    average_precision_score, accuracy_score)
from preprocess import build_features

df = pd.read_csv("../data/upi_100k_ultra_realistic.csv")
df = build_features(df)

drop_cols = ["fraud_flag", "transaction_id",
             "timestamp", "sender_id", "receiver_id"]

X = df.drop(columns=[c for c in drop_cols if c in df.columns])
y = df["fraud_flag"]
\end{lstlisting}

\subsection{Data Splitting}

\begin{lstlisting}[caption={Stratified three-way data split}, label={lst:split}]
# First split: 80% temp, 20% test
X_temp, X_test, y_temp, y_test = train_test_split(
    X, y, test_size=0.2, stratify=y, random_state=42
)

# Second split: 80% train, 20% validation (of temp set)
X_train, X_val, y_train, y_val = train_test_split(
    X_temp, y_temp, test_size=0.2,
    stratify=y_temp, random_state=42
)
\end{lstlisting}

\subsection{Class Imbalance Handling and Model Training}

\begin{lstlisting}[caption={Imbalance weighting and XGBoost training}, label={lst:training}]
# Compute class weight ratio
scale_pos_weight = (
    (len(y_train) - sum(y_train)) / sum(y_train)
)

model = xgb.XGBClassifier(
    n_estimators=800,
    max_depth=10,
    learning_rate=0.02,
    min_child_weight=5,
    gamma=1,
    subsample=0.9,
    colsample_bytree=0.9,
    reg_alpha=0.5,
    reg_lambda=1,
    scale_pos_weight=scale_pos_weight,
    eval_metric="aucpr",
    random_state=42,
    n_jobs=-1
)

model.fit(X_train, y_train)
\end{lstlisting}

\subsection{Threshold Optimization}

\begin{lstlisting}[caption={Validation-set cost-based threshold finding}, label={lst:thresh}]
C_FN = 5000   # Rs. cost per missed fraud
C_FP = 200    # Rs. cost per false positive

val_proba = model.predict_proba(X_val)[:, 1]

def find_best_threshold(y_true, y_prob):
    best_threshold = 0.5
    lowest_loss = float("inf")

    for t in np.arange(0.01, 0.5, 0.01):
        y_pred_temp = (y_prob >= t).astype(int)
        tn, fp, fn, tp = confusion_matrix(
            y_true, y_pred_temp).ravel()
        loss = (fn * C_FN) + (fp * C_FP)

        if loss < lowest_loss:
            lowest_loss = loss
            best_threshold = t

    return best_threshold, lowest_loss

best_threshold, val_loss = find_best_threshold(
    y_val, val_proba)
\end{lstlisting}

\subsection{Model Persistence}

\begin{lstlisting}[caption={Saving trained model artifacts}, label={lst:save}]
pickle.dump(model,
    open("../model/fraud_model.pkl", "wb"))
pickle.dump(X.columns,
    open("../model/feature_columns.pkl", "wb"))
pickle.dump(best_threshold,
    open("../model/threshold.pkl", "wb"))

print("Model, Features, Threshold Saved Successfully")
\end{lstlisting}

\section{Implementation of the FastAPI Endpoint}

\begin{lstlisting}[caption={FastAPI application setup and model loading}, label={lst:fastapisetup}]
from fastapi import FastAPI
import pickle
import pandas as pd
import os, sys

BASE_DIR = os.path.dirname(os.path.dirname(
    os.path.abspath(__file__)))
sys.path.append(BASE_DIR)
from src.preprocess import build_features

app = FastAPI()

# Load model artifacts at application startup
model     = pickle.load(open(
    os.path.join(BASE_DIR, "model", "fraud_model.pkl"), "rb"))
features  = pickle.load(open(
    os.path.join(BASE_DIR, "model", "feature_columns.pkl"), "rb"))
threshold_path = os.path.join(BASE_DIR, "model", "threshold.txt")
threshold = float(open(threshold_path).read()) \
    if os.path.exists(threshold_path) else 0.5
\end{lstlisting}

Loading model artifacts at startup (rather than on each request) is essential for performance. Loading a pickle file takes tens of milliseconds; doing it once at startup means that prediction latency is almost entirely determined by the model inference time, not I/O.

\section{Implementation of the Streamlit Dashboard}

\subsection{Session State Management}

The dashboard uses Streamlit's session state mechanism to persist prediction results between user interactions:

\begin{lstlisting}[caption={Session state initialization}, label={lst:sessionstate}]
for k in ["prob", "decision", "contributors"]:
    if k not in st.session_state:
        st.session_state[k] = None
\end{lstlisting}

Without this initialization, accessing \texttt{st.session\_state.prob} before any prediction has been run would raise a \texttt{KeyError}. Initializing all state keys to \texttt{None} at startup is the standard Streamlit pattern for handling this.

\subsection{Risk Gauge Visualization}

\begin{lstlisting}[caption={Plotly gauge chart for fraud risk display}, label={lst:gauge}]
import plotly.graph_objects as go

fig = go.Figure(go.Indicator(
    mode="gauge+number",
    value=pct,  # Fraud probability as percentage
    title={'text': "Fraud Risk (%)"},
    gauge={
        'axis': {'range': [0, 100]},
        'steps': [
            {'range': [0,  25], 'color': '#00ff9d'},
            {'range': [25, 50], 'color': '#ffaa00'},
            {'range': [50, 75], 'color': '#ff5733'},
            {'range': [75,100], 'color': '#ff0033'}
        ]
    }
))
st.plotly_chart(fig, use_container_width=True)
\end{lstlisting}

The gauge has four color zones: green (0-25\%, low risk), amber (25-50\%, moderate risk), orange-red (50-75\%, high risk), and deep red (75-100\%, very high risk). This color coding allows an analyst to assess risk at a glance without reading the exact number.

\section{Running the System}

\subsection{Training the Model}

\begin{lstlisting}[language=bash, caption={Commands to train the model}]
cd src
python train.py
\end{lstlisting}

\subsection{Optimizing the Threshold}

\begin{lstlisting}[language=bash, caption={Commands to optimize the threshold}]
cd src
python optimize_threshold.py
\end{lstlisting}

\subsection{Starting the API Server}

\begin{lstlisting}[language=bash, caption={Starting the FastAPI server}]
cd upi-guard
python -m uvicorn api.main:app --reload
# API available at: http://127.0.0.1:8000
# Swagger docs at: http://127.0.0.1:8000/docs
\end{lstlisting}

\subsection{Starting the Dashboard}

\begin{lstlisting}[language=bash, caption={Starting the Streamlit dashboard}]
# In a separate terminal
python -m streamlit run ui/dashboard.py
# Dashboard available at: http://localhost:8501
\end{lstlisting}

\subsection{Testing the API with a Sample Request}

\begin{lstlisting}[language=Python, caption={Sample API request using Python requests library}]
import requests

payload = {
    "transaction_id": "TXN_TEST_001",
    "timestamp": "2024-06-15 02:30:00",
    "sender_id": "USER_12345",
    "receiver_id": "USER_99999",
    "amount": 48000,
    "transaction_type": "P2P",
    "merchant_category": "Retail",
    "sender_state": "Maharashtra",
    "receiver_state": "Delhi",
    "sender_bank": "HDFC",
    "receiver_bank": "SBI",
    "device_type": "Android",
    "network_type": "4G",
    "account_age_days": 15,
    "txn_velocity_1h": 6
}

response = requests.post(
    "http://127.0.0.1:8000/predict",
    json=payload
)

print(response.json())
# Expected output (high-risk transaction):
# {"fraud_probability": 0.921, "decision": "Fraud"}
\end{lstlisting}
